\documentclass[11pt]{beamer}

%%%%%%%% tema e cor %%%%%%%%
\mode<presentation> {
\usetheme{Madrid}
%\usecolortheme{albatross}
}

\usepackage[utf8]{inputenc}
\usepackage{graphicx} 
\usepackage{booktabs}
\usepackage{subcaption}


\institute[HU-Berlin] 
{
\includegraphics[width=1.8cm]{img/logo-hu.png}
}

%%%%%%%% titulo e subtitulo %%%%%%%%
\title[Maximum Cardinality Matching]{Maximum Cardinality Matching Problem} 

%%%%%%%% nome dos autores %%%%%%%%
\author [Bojikian, Witkowski, Vogel]{Narek Bojikian, Piotr Witkowski, Martin Vogel} 

\begin{document}

\begin{frame}
\titlepage 
\end{frame}

%%%%%%%% slides %%%%%%%%
\section{Introduction} 
\begin{frame}{Assignment Problem}
\begin{figure}
    \centering
    \includegraphics[height=5cm]{img/introduction/assignment.png}
\end{figure}
\end{frame}

\begin{frame}{Definitions}
    \begin{itemize}[<+->]
    \item Matching
    \item \textbf{Perfect} Matching
    \item \textbf{Bipartite} Matching
    \item \textbf{Maximum} Matching
    \end{itemize}
  
  %add different figure for perfect matching and bipartite matching  
\only<1>{
	\begin{figure}
	\begin{subfigure}{.5\textwidth}
	\centering
	\includegraphics[width=.4\linewidth]{img/introduction/graph01.eps}	
	\end{subfigure}%
	\begin{subfigure}{.5\textwidth}
	\centering
	\includegraphics[width=.4\linewidth]{img/introduction/graph02.eps}
	\end{subfigure}
	\end{figure}
}
\only<2>{
	\begin{figure}
	    \centering
	    \includegraphics[width=0.3\linewidth]{img/introduction/perfectmatching.eps}
	\end{figure}
}
\only<3>{
	\begin{figure}
	\begin{subfigure}{.5\textwidth}
	\centering
	\includegraphics[width=.4\linewidth]{img/introduction/bipar.eps}
	\end{subfigure}%
	\begin{subfigure}{.5\textwidth}
	\centering
	\includegraphics[width=.4\linewidth]{img/introduction/biparmatch.eps}
	\end{subfigure}
	\end{figure}
}
\only<4>{
	\begin{figure}
	\begin{subfigure}{.5\textwidth}
	\centering
	\includegraphics[width=.4\linewidth]{img/introduction/maximalmatch.eps}
	\end{subfigure}%
	\begin{subfigure}{.5\textwidth}
	\centering
	\includegraphics[width=.4\linewidth]{img/introduction/maximummatching.eps}
	\end{subfigure}
	\end{figure}
}
    
\begin{block}{Definition}
    \only<1>{A set of non overlapping edges.}
    \only<2>{Matching of size $\frac{|V\left(G\right)|}{2}$}
    \only<3>{Matching on a bipartite graph.}
    \only<4>{Matching of the maximum size.}
\end{block} 
\end{frame}

\begin{frame}{Maximum Matching vs Maximal Matching}
	\parbox[c][.7\textheight][c]{\textwidth}{%
		\only<1>{
			\begin{center}
				\includegraphics[width=0.15\linewidth]{img/random/maximum-maximal-0.eps}
			\end{center}
		}
		\only<2>{
			\begin{center}
				\includegraphics[width=0.15\linewidth]{img/random/maximum-maximal-1.eps}
			\end{center}
		}
		\only<3->{
			\begin{center}
				\includegraphics[width=0.15\linewidth]{img/random/maximum-maximal-2.eps}
			\end{center}
		}
		\uncover<4>{
			\begin{center}
				\includegraphics[width=0.15\linewidth]{img/random/maximum-maximal-3.eps}
			\end{center}
		}
	}
\end{frame}


\section{Justificativa}
\input{files/justificativa.tex}

\section{Objetivos}
\input{files/objetivos.tex}

\section{Fundamentação Teórica}
\input{files/fundamentacao.tex}

\section{Metodologia}
\input{files/metodologia.tex}

\section{Resultados}
\input{files/resultados.tex}

\section{Conclusão}
\input{files/conclusao.tex}

%%%%%%%% agradecimentos %%%%%%%%
\begin{frame}{Agradecimentos}
    \large{Agradeço a fulano, ciclano e beltrano que apoiaram o desenvolvimento dessa pesquisa.}
\end{frame}
%------------------------------------------------


%------------------------------------------------
%------------------------------------------------
%%%%%%%% referencias %%%%%%%%
\nocite{*}
\begin{frame}[allowframebreaks]{Referências}
\bibliographystyle{unsrt}
\bibliography{ref.bib}
\end{frame}


%%%%%%%% repete primeiro slide %%%%%%%%
\begin{frame}
\titlepage 
\end{frame}


\end{document}